\begin{table*}[t]
\centering\small
\begin{tabular}{llrrrr}
\toprule
Panel & Method & Setup & Cal. & Cand. test & Accounted total \\
\midrule
RewardBench 2 & Reference & 48244 & 0 & 285.0 & 48529.0 \\
RewardBench 2 & Full batch & 48244 & 0 & 16100.0 & 64344.0 \\
RewardBench 2 & CRC & 48244 & 32144 & 16100.0 & 96488.0 \\
RewardBench 2 & SCoRE & 48244 & 32144 & 0.0 & 80388.0 \\
WMT23, 13 streams & Reference & 74737 & 0 & 523.2 & 75260.2 \\
WMT23, 13 streams & Full batch & 74737 & 0 & 41327.0 & 116064.0 \\
WMT23, 13 streams & CRC & 74737 & 57720 & 41327.0 & 173784.0 \\
WMT23, 13 streams & SCoRE & 74737 & 57720 & 0.0 & 132457.0 \\
WMT23, 12 streams & Reference & 68988 & 0 & 43.3 & 69031.3 \\
WMT23, 12 streams & Full batch & 68988 & 0 & 38148.0 & 107136.0 \\
WMT23, 12 streams & CRC & 68988 & 53280 & 38148.0 & 160416.0 \\
WMT23, 12 streams & SCoRE & 68988 & 53280 & 0.0 & 122268.0 \\
\bottomrule
\end{tabular}
\caption{Full batch denotes the full-information benchmark, not an online acquisition policy. At $q=.05$, median pooling, $\alpha=.10$: logical score-entry accounting. Setup includes pool/candidate reference and pool test entries; Cal.\ includes pool/candidate calibration entries; Cand. test excludes base-test scores already counted in Setup. Accounted total sums these three columns after panel construction; it omits preprocessing reads of the full panel. Values average 40 splits and do not measure runtime. Endpoint zero-query certification can omit candidate-reference scoring; no setup-optimality claim is made.}
\label{tab:review-workload}
\end{table*}
